% =====================================================================
%  Conferenze, Zagarolo (On The Road, III Conferenza internazionale sul
%  Game Design), giorno 1 — Intervista 1
%  Il marketing di un publisher indie: che cosa fa, che cosa cerca
%  Fonte: trascrizione automatica "Track 9". Intervista in italiano a una
%  persona del reparto marketing del publisher britannico Fireshine
%  Games, condotta da un moderatore dell'organizzazione, con domande dal
%  pubblico. La registrazione comincia durante le presentazioni.
%  Etichette con prefisso "c2-1:". Nessun pacchetto aggiuntivo.
% =====================================================================

\chapter[Il marketing di un publisher indie]{Il marketing di un publisher
indie: che cosa fa, che cosa cerca}
\label{cap:c2-1-marketing-publisher}

Il primo giorno di \emph{On The Road} a Zagarolo si apre, nella parte registrata,
con un'intervista dedicata al \textbf{marketing dal punto di vista di un publisher
indipendente}. L'ospite lavora nel marketing di \textbf{Fireshine
Games}\footnote{Trascritto anche ``Fire Chain Games'', ``Farsha'', ``Firestar'',
``Feierstein'', ``Fashion''. Fireshine Games è un publisher indipendente
britannico.}, publisher britannico, ed è la prima di due voci della stessa
azienda: subito dopo interverrà un producer (capitolo~\ref{cap:c2-2-producer}).
Il moderatore, dell'organizzazione dell'evento, avverte che ``siccome non ci sono
telecamere'' si può parlare liberamente, e invita il pubblico a fare le domande
che molti studi di solito non fanno ``per essere abbastanza educati''.

% ---------------------------------------------------------------------
\section{L'ospite}
\label{sec:c2-1-ospite}

L'ospite\footnote{Il nome è trascritto in modi diversi (``Aerto'', ``Larry'',
``l'AERT'', ``l'art''), tutti probabilmente storpiature di un nome albanese
breve; non è stato possibile ricostruirlo. Nel capitolo si parla semplicemente
dell'``ospite''.} si presenta da sé: il nome ``strano'' si deve al fatto di essere
nata in \textbf{Albania} e arrivata in Italia a un anno. Ha vissuto in Italia fino
a dodici anni fa, poi è partita per l'America (Canada, East Coast, \textbf{New
York}). Innamorata di New York ma senza un visto facile, ha scelto ``la cosa più
vicina'': \textbf{Londra}, ``la New York europea'', prima della Brexit, quando
bastava presentarsi e chiedere un lavoro.

Dopo un po' di gavetta ha studiato alla \textbf{National Film and Television
School} (NFTS), \textbf{marketing per film, TV e videogiochi}. Ha capito presto che
il cinema non faceva per lei (``un'industria molto avversa al rischio'') e, da
giocatrice, ha provato i videogiochi:
\begin{itemize}
  \item ha cominciato come \textbf{stagista} presso un publisher indie,
        \textbf{Curve Games};
  \item dopo due anni ha \textbf{perso il lavoro} (``capita, quando le cose non
        vanno bene'');
  \item grazie a un amico conosciuto all'NFTS ha trovato subito un posto in
        \textbf{Fireshine}, dove lavora da tre anni e dove pensa di restare
        ``almeno per altri tre''.
\end{itemize}

% ---------------------------------------------------------------------
\section{Che cosa fa un publisher nel marketing}
\label{sec:c2-1-publisher}

Molti studi affidano tutto il marketing al publisher, ma si chiedono che cosa
faccia esattamente e quanto budget usi.

Gli studi indie, risponde l'ospite, cercano un publisher per due cose:
\textbf{i soldi} per fare il gioco e \textbf{il marketing}. Ma un publisher porta
molto di più: un \textbf{sapere quasi nascosto} all'industria. Un publisher indie
lancia \textbf{almeno sei giochi all'anno}; da ogni lancio impara molto, è sempre
in contatto con l'industria, capisce i \textbf{trend} del marketing ed è in
contatto con \textbf{partner globali}: pubblicità in Giappone, Cina, Sud America,
Nord America ed Europa.

Il contributo più grosso è \textbf{aprire porte} che chi si autopubblica non ha, a
meno di essere uno studio con un grande successo alle spalle (che diventa ``quasi
un publisher''):
\begin{itemize}
  \item gli \textbf{showcase digitali}\footnote{I nomi sono storpiati (``AAA o il
        Jack Kelly Showcase, Indy Direct''); si riconoscono il riferimento agli
        showcase indipendenti e a quelli di PlayStation e Xbox.}, compresi quelli
        di PlayStation e Xbox;
  \item un team che rappresenta il gioco \textbf{a livello globale}, per esempio
        per una versione mobile o per trovare un partner in Giappone o in Cina.
\end{itemize}

\subsection{Il marketing come ``cattivo'' dell'industria}

Una cosa che molti sottovalutano: il marketing è considerato \textbf{il cattivo
dell'industria}. Se uno sviluppatore fa qualcosa di sbagliato, molti dicono
``sicuramente è stato il marketing a dirgli di farlo''. Ma ``nel 90\% dei casi''
è stato lo sviluppatore, o comunque un'idea che non funzionava. Almeno nell'indie,
il publisher cerca di dare allo studio \textbf{totale libertà}, nei limiti del
budget.

\subsection{Quanto conta il marketing nella scelta dei giochi}

Il marketing influenza la decisione di prendere un gioco? Sì, ma
\textbf{indirettamente}. In Fireshine c'è un \textbf{team di acquisizione} guidato
da una persona (``Jasper'') che ha pieno potere sulla strategia e sul portafoglio.
Il marketing gli spiega costantemente che cosa funziona e che cosa no sul mercato,
e quando un gioco interessante entra nel suo radar gli viene chiesto un parere.

In Fireshine ``tutti sono giocatori'', anche chi lavora nell'amministrazione. Così,
quando arriva un gioco interessante, lo provano tutti, e \textbf{se scatta
l'innamoramento ``crollano gli schemi''}: lo vogliono fare anche se dal punto di
vista del marketing non ha senso. Quando invece non tutti si innamorano, si fanno
i conti:
\begin{itemize}
  \item funziona dal punto di vista del marketing?
  \item il budget ha senso?
  \item lo sviluppatore può portare il progetto a termine? Ha il curriculum per
        farlo?
\end{itemize}
Il marketing risponde a una di queste domande, cercando di essere obiettivo, ``ma
alla fine nessuno sa veramente che cosa funziona'': è sempre un rischio.

% ---------------------------------------------------------------------
\section{Consigli agli sviluppatori italiani}
\label{sec:c2-1-consigli}

\subsection{Un'idea che si racconta in trenta secondi}

Il consiglio principale, soprattutto per i team piccoli: avere \textbf{un'idea
chiara}, che si possa raccontare \textbf{in 30 secondi}, che possa passare
\textbf{in un post di Instagram}. Più l'idea è chiara e cristallina, più è facile
venderla: non solo ai giocatori, ma a chi deve darti ``mezzo milione'' per fare
il gioco. I publisher hanno molto lavoro: una mail si legge in due minuti, un
\emph{deck} di dieci pagine in cinque. Se l'idea non è forte e trasmessa
chiaramente, è difficile arrivare al secondo passaggio.

Altri consigli sul pitch:
\begin{itemize}
  \item \textbf{pitch troppo lunghi non funzionano}: corto, semplice, diretto;
  \item non preoccuparsi troppo del budget: se il publisher è quello giusto, non
        sarà un problema;
  \item ma \textbf{né troppo né troppo poco}: i publisher sono esperti. Chi chiede
        troppo riceve un no; chi chiede troppo poco rivela di non avere
        esperienza, cioè di non sapere che cosa vuole fare. Bisogna saper
        \textbf{motivare la cifra} richiesta.
\end{itemize}

\subsection{Il costo del lavoro e, soprattutto, il tempo}

Il moderatore osserva che in Italia il costo del lavoro è molto più basso che nel
Regno Unito: lì 3.000--4.000 sterline al mese per una risorsa sono normali, qui
3.000--3.500 euro sono ``quasi impensabili''. Il publisher ne tiene conto, o pensa
che un budget italiano sia troppo basso?

Se ne tiene conto, risponde l'ospite, ma conta di più \textbf{il tempo}. Un team di
cinque persone che impiega \textbf{due anni} è un costo enorme, e non si sa se tra
due anni l'idea avrà ancora senso o se il team si sarà perso. L'ideale è qualcosa
che si possa eseguire in modo efficiente in \textbf{8--12 mesi}. Chi ha già avuto
due o tre successi, o anche uno solo, e ha dimostrato di conoscere il processo, può
discutere anche di budget molto più alti. Che il team voglia essere pagato 3.000 o
5.000 al mese ``va bene'', ma ci sono anche i costi di QA, localizzazione,
marketing e tanto altro, più un team di marketing e business che lavora intorno al
progetto. Quindi: \textbf{come primo gioco, fate qualcosa che sia fattibile in
6--12 mesi}. Il moderatore sottolinea quanto sia utile l'informazione, perché in
Italia vede moltissimi pitch da \textbf{18--24 mesi}.

\subsection{L'ossessione delle wishlist}

A un publisher interessa che il gioco sia già su Steam, o è meglio che non sia
annunciato? Quante wishlist servono?

``Un milione sarebbe bello.'' Le wishlist contano: un primo gioco con 100.000 o
200.000 wishlist avrà ``tutti i publisher ai suoi piedi'' e un enorme potere
contrattuale; con zero, le cose cambiano, e se il gioco non è nemmeno annunciato è
solo un concept. Eppure \textbf{i progetti più importanti di Fireshine avevano zero
wishlist} al momento del via libera:
\begin{itemize}
  \item \emph{Far Far West}\footnote{Trascritto anche ``Fafa West'', ``FIFA
        West'', ``Papa West''.} non era neanche annunciato, e molti publisher
        avevano detto no;
  \item un altro progetto era ``un'idea'' molto particolare, rifiutata da molti
        publisher\footnote{La trascrizione dice ``Dance Attack era un'idea. C'è Tony
        Hawk e un treno'': non è stato possibile identificare il gioco.};
  \item un \textbf{3D action platform}\footnote{Trascritto ``Da Spade'' e
        ``Daspade''; titolo non verificabile.}, genere ``non fortissimo al giorno
        d'oggi''.
\end{itemize}
Probabilmente il successo è venuto dal fatto che quegli sviluppatori avevano
\textbf{già fatto un gioco} (autopubblicato o in aziende che pubblicavano giochi):
avevano \textbf{un curriculum importante}, un'idea chiara e hanno dimostrato di
poter portare a termine il progetto. Questo conta più delle wishlist e dei KPI.
Regola pratica: \textbf{se hai già annunciato il gioco, devi avere le wishlist};
se non l'hai annunciato, hai una possibilità con un pitch forte o un background
importante.

\subsection{Come arrivare a centomila}

Per un \emph{solo dev} o chi si autopubblica funzionano molto gli \textbf{short}
su TikTok, Instagram e YouTube. Non costano nulla a parte il tempo, e ogni video
\textbf{è un pitch}: se non vedi alcuna crescita né interesse, probabilmente
l'idea è prematura o non hai le competenze per eseguirla; se vedi anche solo una
crescita normale, c'è qualcosa. L'interesse attira publisher, fondi e
collaboratori. Bisogna saper comunicare il gioco su queste piattaforme, preparare
un \textbf{pitch professionale} e \textbf{insistere}: ``vi si chiuderanno tante
porte, ma un giorno una si apre''. Gli eventi in presenza aiutano a conoscere
persone e sviluppatori, ognuno con la propria storia: ``non c'è un format''.
L'ospite dice di aver visto, in sei anni, ``gente che era nessuno'' fare un gioco
e avere successo. Oggi vanno forte TikTok e Instagram; domani sarà
qualcos'altro.

% ---------------------------------------------------------------------
\section{Le domande del pubblico}
\label{sec:c2-1-domande}

\subsection{Vendere il gioco o mostrarlo?}

Uno studente di game design chiede: il videogioco è sia prodotto sia opera
espressiva; il marketing deve \emph{vendere} il gioco o \emph{mostrarlo}?

La regola del reparto è chiara: \textbf{``noi non vendiamo il gioco''}. Si cerca di
portarlo davanti a quante più persone possibile, ed \textbf{è il gioco a trovare
il suo pubblico}. Nel 2026 non si può forzare un prodotto sulle persone, a meno di
essere Apple o Netflix. Per questo i giochi indie di maggior successo sono quelli
che hanno avuto \textbf{una demo in anticipo}, \textbf{playtest}, influencer che li
hanno giocati prima: hanno \textbf{garantito la qualità prima di chiedere i
soldi}. Per \emph{Far Far West} Fireshine ha fatto due playtest aperti e una demo:
\textbf{mezzo milione di giocatori} in tre test, con commenti del tipo ``non ci
posso credere che sia gratis, potreste lanciarlo adesso''. Lo sviluppatore, da
professionista, voleva fare ancora di più; ma a quel punto il publisher sapeva che
il gioco avrebbe venduto moltissimo. Molti giochi, invece, il loro pubblico non lo
trovano mai, ``e ovviamente è colpa nostra''\footnote{La trascrizione ripete la
frase decine di volte per un errore meccanico.}.

\subsection{L'IA e i cloni}

Con l'IA è più facile creare prototipi e metterli sul mercato; quando un gioco ha
successo escono subito molte copie. Se la meccanica non è più un vantaggio, su che
cosa puntare?

Per un \emph{solo dev} è ``veramente difficile'': un'idea pazzesca può essere
copiata da qualcuno con un marchio più forte. I giochi di Fireshine non sono facili
da copiare: dipendono molto dal \textbf{talento} di chi ci lavora; alcuni sono idee
straordinarie, altri idee non pazzesche ma \textbf{eseguite in modo eccezionale}.
Le contromisure:
\begin{itemize}
  \item \textbf{campagne più corte}: dove prima si facevano 12 mesi di marketing,
        ora si pensa a 3--4 mesi, per non dare ai cloni il tempo di uscire; e se
        escono, è troppo tardi;
  \item il \textbf{marketing} è difficile da copiare: chi clona non ha il potere di
        fare mezzo milione di wishlist in un mese;
  \item come publisher, Fireshine ha un \textbf{account manager in Valve} a cui
        segnalare un clone perché venga rimosso; un \emph{solo dev} spesso non
        può difendersi.
\end{itemize}
Diverso è il caso dei \textbf{\emph{friend slop}}\footnote{Trascritto ``friend
slot'': i giochi cooperativi caotici ed economici da giocare con gli amici.}: un
genere che molti davano per finito e invece ``è qui per restare'', facile da
copiare perché va virale; su Steam il giocatore può trovare il clone invece
dell'originale. ``L'industria si adatterà.'' Paradossalmente, oggi i dev ancora più
spaventati cercano un publisher, e questo è un vantaggio per i publisher.

\subsection{Il rischio delle dirette}

Le dirette degli influencer non rischiano di far ``consumare'' il gioco a chi si
accontenta di guardarlo? Non si saprà mai se chi guarda su Twitch comprerà, dice
l'ospite. L'obiettivo del marketing è creare la sensazione di \textbf{non potersi
perdere un momento culturale}: se tutti giocano a \emph{GTA}, non vuoi solo
guardare, vuoi entrare. Con \emph{Far Far West} molti guardavano il proprio
streamer preferito, ma molti hanno detto ``sembra divertente, voglio giocarci con i
miei amici''. Chi guarda solo, probabilmente non l'avrebbe comprato comunque.

E il marketing non è solo promozione: comprende decidere il \textbf{prezzo} (che
varia tra regioni), scegliere \textbf{quali content creator} far giocare,
preparare gli \textbf{asset video} che convincono all'acquisto, costruire la
\textbf{pagina Steam} (``c'è uno studio dietro''). Spesso un gioco arriva al lancio
\textbf{avendo già recuperato i costi} grazie a un accordo con \textbf{Game Pass} o
a un \textbf{co-publishing} in Cina o in Giappone: il content creator è solo una
piccola parte del lavoro.

\subsection{Prima la storia o il gameplay?}

Un filmmaker chiede se un gioco debba partire dalla storia o dal gameplay.
Dipende dal genere: molti giochi sono \textbf{sistemici}, e il sistema conta più
della storia. Nell'indie Fireshine fa giochi da \textbf{300.000 a un milione e
mezzo} di budget. La cosa più importante è che il gioco sia \textbf{divertente},
che tu voglia dire a tutti ``questo gioco è fantastico''. Per il platform 3D la
storia era importante e c'erano due \emph{writer}; ma due titoli che insieme hanno
venduto \textbf{6 milioni di copie} (\emph{Far Far West} e un altro\footnote{
Trascritto ``Call Keeper''; il titolo non è verificabile.}) di fatto \textbf{non
hanno una storia scritta}: hanno un \textbf{mondo} costruito bene, che racconta
comunque e attira i giocatori.

Lo stesso spettatore, a cui hanno detto che i suoi film ``sembrano videogiochi'',
chiede come un regista e sceneggiatore possa avvicinarsi al settore. Il consiglio:
contattare studi che fanno giochi in cui lo \textbf{storytelling} è una colonna
portante (cita \emph{Dispatch}, ``pazzesco'', in cui la storia è tutto); nei giochi
di Fireshine la storia è una colonna ma non l'unica, e a volte non lo è affatto.

\subsection{Che cosa fa spiccare un pitch}

Che cosa colpisce in un pitch, e come si interviene su un'idea che piace?

\begin{itemize}
  \item Fireshine \textbf{non cambia le idee}: fa \textbf{domande scomode}. Il team
        di scouting è ``veramente bravo'': se non hai le risposte, stai sviluppando
        un'idea che non sei pronto a sviluppare, e devi tornare al pitch. ``Non ti
        diremo mai che cosa fare'', altrimenti non sarebbero un publisher indie.
  \item Dal punto di vista \textbf{visivo}, conta ciò che va su Steam e
        \textbf{l'identità del publisher}: a un publisher molto ricercato e indie
        non si porta un gioco che sembra \emph{Fortnite}; a Devolver si porta
        qualcosa che sembra un gioco Devolver; a Fireshine qualcosa che potrebbe
        essere un gioco Fireshine.
  \item Un'idea chiara: lo scout deve leggere la prima frase e riuscire a
        \textbf{visualizzarla}, e poi voler aprire la presentazione.
  \item E la fortuna: ``molto importante''.
\end{itemize}
Anche quando un progetto è buono ma non adatto, Fireshine dà \textbf{feedback} e
suggerisce altri publisher, oppure dice onestamente che forse \textbf{non serve un
publisher} e conviene autopubblicarsi, perché con il publisher lo studio non
guadagnerebbe nulla e il publisher non muoverebbe il proprio business. Altre volte
dice: ``torna tra sei mesi con qualcosa che possiamo giocare''.

\subsection{Come tutelarsi sul budget di marketing}

Uno sviluppatore chiede come assicurarsi che il budget di marketing venga speso
bene: è capitato a più studi di vedere budget grossi spesi male (``30.000 euro per
un trailer'') o fatti circolare internamente dal publisher.

L'ospite ammette che ``molti publisher fanno quello''; Fireshine no. Pur essendo
``una che spende tanto'' in marketing, non ha mai speso 30.000 euro per un trailer:
i loro trailer costano \textbf{2.000--3.000 euro}, a volte nulla, grazie a video
editor interni. Il modo di tutelarsi è scegliere un \textbf{publisher della
dimensione giusta}: uno grande, con una struttura costosa, spenderà molto perché è
abituato a farlo; uno della dimensione giusta opererà con budget minori. Su un
accordo da \textbf{un milione di sterline}, Fireshine spende circa
\textbf{250.000} in marketing. Ma la spesa è \textbf{flessibile}: se il gioco non
fa wishlist, non si spende, e si dice allo studio che l'idea non sta funzionando;
se fa mezzo milione di wishlist, si spende ancora di più. Con i publisher seri si
può avere accesso al budget, vedere come viene speso e chiarire prima dell'accordo
che cosa non si vuole.

\subsection{Che cosa pensa Fireshine degli studi italiani?}

L'ultima domanda, ``scomoda'', è del moderatore: Fireshine viene spesso in Italia;
che cosa pensa davvero degli studi italiani?

Gli scout \textbf{non guardano la nazionalità}, ``né che cosa mangi o dove vivi'':
guardano il gioco, l'esperienza, il budget, l'adattamento al portafoglio. Detto
questo, secondo l'ospite \textbf{l'industria italiana è indietro}: gli svedesi e
i nordeuropei sono ``anni luce avanti'', e la Spagna, pur con economia e cultura
simili, ``fa benissimo''. Speculando, il problema potrebbe essere una
\textbf{società vecchia}, in cui i giovani non trovano spazio (``sei troppo
giovane, devi fare esperienza''). Il consiglio: è un'industria globale,
\textbf{viaggiate}, potete crescere anche fuori dall'Italia. E all'Italia
suggerisce di \textbf{investire}: è un'industria in cui ``con 100.000 euro si
possono fare 20 milioni'', e ha visto ``ragazzini fare 20 milioni in un anno'', con
tutte le tasse che ne derivano. Il moderatore aggiunge che un acceleratore tedesco
incontrato la settimana prima gli ha raccontato di circa \textbf{50 milioni di
euro} investiti nell'industria nella sola regione di Francoforte. L'ospite chiude
con un esempio: se 200 sviluppatori arrivano con un'idea e, alla domanda ``quanti
soldi vi servono?'', rispondono ``nessuno, il governo mi ha già finanziato'', il
publisher ne prende un paio, e se uno fa un sacco di soldi ``il governo fa tanti
soldi''.

% ---------------------------------------------------------------------
\section*{Punti chiave del capitolo}
\addcontentsline{toc}{section}{Punti chiave del capitolo}

\begin{itemize}
  \item Un publisher indie porta soldi e marketing, ma soprattutto
        \textbf{conoscenza} (sei o più lanci all'anno), accesso agli
        \textbf{showcase} e una rete di \textbf{partner globali}.
  \item Il marketing influenza le acquisizioni indirettamente; un gioco di cui
        tutti si innamorano ``fa crollare gli schemi''; per gli altri contano
        mercato, budget e \textbf{curriculum} del team.
  \item Pitch: un'idea raccontabile \textbf{in 30 secondi}, corto e diretto,
        budget né troppo alto né troppo basso ma \textbf{motivato}.
  \item Conta più \textbf{il tempo} che il costo del lavoro: per un primo gioco,
        6--12 mesi, non 18--24.
  \item Le wishlist danno forza contrattuale, ma i maggiori successi di
        Fireshine sono partiti da zero, con team che avevano già pubblicato. Se il
        gioco è annunciato, le wishlist servono.
  \item Gli short (TikTok, Instagram, YouTube) sono pitch gratuiti che misurano
        l'interesse. ``Noi non vendiamo il gioco'': demo e playtest garantiscono la
        qualità prima di chiedere soldi.
  \item Contro i cloni: campagne più corte (3--4 mesi), forza del marketing,
        segnalazioni a Valve.
  \item Il marketing è anche prezzo regionale, scelta dei creator, asset, pagina
        Steam; spesso i costi sono coperti prima del lancio da accordi (Game Pass,
        co-publishing in Asia).
  \item Per tutelare il budget: un publisher \textbf{della dimensione giusta};
        spesa flessibile (circa il 25\% di un accordo in marketing, solo se il
        gioco risponde).
  \item L'industria italiana è indietro; servono investimenti e la disponibilità
        a muoversi in un'industria globale.
\end{itemize}
